% ============================================================
% Chapter 25: Retrospective — The Legacy of Hilbert's Problems
% ============================================================

\chapter*{Retrospective: The Legacy of Hilbert's Problems}
\addcontentsline{toc}{chapter}{Retrospective: The Legacy of Hilbert's Problems}

We have now journeyed through all twenty-three of Hilbert's problems. Each chapter has told its own story---of a question posed, of tools forged, of answers found or still sought. But a book about twenty-three problems is also a book about something larger: the nature of mathematical ambition, the arc of mathematical progress, and the curious fact that the most significant questions of one century become the foundational theorems of the next.

This final chapter steps back from the individual problems to survey the landscape as a whole. What did Hilbert get right? What did he get wrong? And what can his list teach us about the enterprise of mathematics itself?

% ============================================================
\section{What Hilbert Got Right}
% ============================================================

The most obvious measure of success is this: Hilbert's problems \emph{worked}. They guided research, created new fields, and led to some of the most significant theorems of the twentieth century. The list was not a random collection of questions---it was a map of the mathematical frontier in 1900, and for the most part, it was a notably accurate one.

\subsection{Problems That Shaped Entire Fields}

Several of Hilbert's problems were not merely questions within existing disciplines---they \emph{created} those disciplines, or at least accelerated their birth.

\paragraph{Problem 1 (Continuum Hypothesis) and Set Theory.} Cantor's work on the infinite was still controversial in 1900. By placing the Continuum Hypothesis first on his list, Hilbert signaled that questions about the structure of the infinite were central to mathematics, not peripheral curiosities. The resolution of CH---via G\"odel's constructible universe and Cohen's forcing---gave birth to modern set theory as a mature discipline. Forcing has become one of the most important techniques in all of mathematical logic, with applications in model theory, recursion theory, and the study of cardinal arithmetic. The independence of CH from ZFC raised philosophical questions about mathematical truth that remain unresolved. None of this would have happened with the urgency it did without Hilbert's signal.

\paragraph{Problem 7 (Irrationality and Transcendence) and Number Theory.} Hilbert's seventh problem---that $a^b$ is transcendental when $a$ is algebraic ($\neq 0,1$) and $b$ is irrational and algebraic---was solved by Gelfond and Schneider in 1934. But the deeper legacy was the development of transcendental number theory as a field. Baker's theorem on linear forms in logarithms (1966), which won the Fields Medal, grew directly from the tradition Hilbert's problem helped establish. The theory of transcendence remains a vibrant area of research, with connections to Diophantine geometry and arithmetic dynamics.

\paragraph{Problem 10 (Decision Problem for Diophantine Equations) and Computability Theory.} Hilbert's tenth problem---is there an algorithm to decide whether a polynomial equation with integer coefficients has integer solutions?---led, through its negative answer by Matiyasevich (1970), to the creation of computability theory as we know it. The connection between Diophantine equations and Turing machines---the MRDP theorem---was one of the most surprising results of the twentieth century. It linked number theory to the theory of computation in a way no one had anticipated. The techniques developed along the way---recursively enumerable sets, Hilbert's Tenth problem in model theory, connections to algebraic geometry over $\mathbb{Z}$---have become standard tools in logic and theoretical computer science.

\paragraph{Problem 6 (Axioms of Physics) and the Foundations of Quantum Mechanics.} Hilbert's sixth problem called for an axiomatic treatment of physics. While this remains unfinished, the program it inspired was significantly productive. von Neumann's axiomatic treatment of quantum mechanics (1932), Kolmogorov's axiomatic foundation of probability theory (1933), and the axiomatic approach to statistical mechanics all bear the mark of Hilbert's vision. The idea that physical theories should be placed on a rigorous mathematical foundation---now taken for granted---was, in 1900, a significant program.

\subsection{Problems That Led to Significant Theorems}

Other problems, while not creating entire fields, led to results of significant depth and elegance.

\paragraph{Problem 3 (Equality of Volumes).} Dehn's invariant, which resolved this problem in 1901, was an important result of combinatorial topology. The idea that geometric questions can be answered by algebraic invariants---that you can prove two things are \emph{not} equivalent by computing a number that distinguishes them---became a paradigm for all of algebraic topology and homological algebra.

\paragraph{Problem 5 (Lie Groups as Manifolds).} The proof that every continuous group is a Lie group (by Cartan, von Neumann, and Pontryagin in the 1930s) was an important result in the theory of Lie groups and Lie algebras. It bridged topology and algebra in a way that presaged the development of modern differential geometry and representation theory.

\paragraph{Problem 9 (Reciprocity Laws).} The generalization of quadratic reciprocity to arbitrary number fields---Hilbert's ninth problem---was addressed through the development of class field theory by Artin, Hasse, and Takagi. Class field theory is one of the significant achievements of twentieth-century mathematics, and its modern avatar, the Langlands program, is arguably the most significant unifying vision in contemporary number theory.

\paragraph{Problem 21 (Monodromy Problem).} The proof of the Riemann--Hilbert correspondence (by Deligne, Grothendieck, and others) was an important result in algebraic geometry. It connected the monodromy group of a differential equation to the geometry of its solutions, bridging analysis and algebra in a significant way. The Riemann--Hilbert correspondence remains a central tool in the study of differential equations and perverse sheaves.

\subsection{The Problems as a Research Program}

What Hilbert got most right was not any individual problem---it was the \emph{idea} of a problem list as a research program. Before Hilbert, mathematicians worked on problems that interested them; after Hilbert, the idea that a community of mathematicians could be organized around a shared set of challenges became natural. The Clay Millennium Problems (discussed below), the Birch and Swinnerton-Dyer conjecture, the Navier--Stokes existence and smoothness problem, the Hodge conjecture---these are all ``Hilbert problems'' in spirit, even if they were not on Hilbert's original list.

% ============================================================
\section{What Hilbert Got Wrong}
% ============================================================

To say that Hilbert's problems ``worked'' is not to say that every problem was well-posed or that every expectation was fulfilled. Some problems were too specific, some rested on assumptions that turned out to be false, and some remain too vague to have a definite answer.

\subsection{Assumptions That Failed}

The most dramatic failure of Hilbert's expectations was Problem 2, the consistency of arithmetic. Hilbert's \emph{program}---his grand plan to prove the consistency of all of mathematics using only ``finitary'' methods---was demolished by G\"odel's incompleteness theorems in 1931. Hilbert believed that a proof of consistency was within reach, and that such a proof would secure the foundations of mathematics once and for all. G\"odel showed that any formal system powerful enough to describe arithmetic, if consistent, cannot prove its own consistency. The dream of a complete and consistent foundation for all of mathematics, achievable by a finite number of finitary steps, died in 1931.

This is not merely a historical footnote. It changed the landscape of mathematical logic permanently. Hilbert's program, in its original form, is impossible. Modified versions (such as proof theory in the sense of Gentzen, or the study of consistency strength) have been significantly productive, but the original vision of certainty through finitary proof was wrong.

\subsection{Problems That Were Too Specific}

Some of Hilbert's problems were so specific that their resolution, while interesting, did not have the broad impact Hilbert anticipated. Problem 13---the solvability of the seventh-degree equation by radicals---was essentially a narrow question in algebra, solved in Hilbert's own time. The result was notable, but it did not open new territory.

Similarly, Problem 18---Keller's conjecture on tessellations of space---was resolved (negatively, by Debrunner in 1965 for $n \geq 6$), but the problem was more of a curiosity than a gateway. It did not lead to a new theory or a new way of thinking.

\subsection{Problems That Remain Too Vague}

Some of Hilbert's problems were more \emph{programs} than problems---invitations to develop a theory rather than specific questions with definite answers. Problem 6 (axioms of physics), Problem 15 (Schubert's enumerative calculus), Problem 16 (topology of algebraic curves), and Problem 23 (calculus of variations) are all of this kind. They have been ``partially solved'' in the sense that significant progress has been made, but they lack the crisp statement that would allow a definitive resolution.

This vagueness is both a strength and a weakness. It is a strength because it allows the problem to evolve as mathematics evolves---the problem can mean different things in different eras, and new techniques can breathe new life into old questions. It is a weakness because it makes it impossible to declare a problem ``solved'' in the way that Problem 3 or Problem 22 has been solved. The boundary between a significant open problem and a research program is often unclear.

\subsection{The Problem of Vagueness: A Case Study}

Consider Problem 16: ``Investigate the topology of algebraic varieties defined by algebraic equations.'' What does it mean to ``investigate'' an entire class of objects? In 1900, algebraic topology was in its infancy. By the 1960s, with the work of Thom, Milnor, Smale, and others, the topology of smooth manifolds was well understood, but the topology of \emph{singular} algebraic varieties---the objects that actually arise from polynomial equations---remained mysterious. Today, the subject is called \emph{singularity theory} and \emph{algebraic topology of varieties}, and it is a significant enterprise. Has Problem 16 been ``solved''? In a sense, yes: we know a great deal. In another sense, no: the investigation continues indefinitely.

This illustrates a general point: Hilbert's problems were not always questions in the narrow sense. They were \emph{visions}---glimpses of a mathematical landscape that Hilbert could see but not fully describe. The value of a vision is not always in its specific content but in the direction it provides.

% ============================================================
\section{The Unsolved Problems Today}
% ============================================================

Of Hilbert's twenty-three problems, the most prominently unsolved is Problem 8, the Goldbach conjecture. Several others are ``partially solved'' in the sense discussed above, but none has the simple, clean statement of the Goldbach conjecture and the significant, stubborn resistance to proof.

\subsection{The Goldbach Conjecture}

\begin{quote}
\textit{Every even integer greater than 2 is the sum of two primes.}
\end{quote}

This conjecture, first stated by Christian Goldbach in a letter to Euler in 1742, has been verified for all even integers up to $4 \times 10^{18}$ by computer search. The heuristic evidence is overwhelming: the number of ways to write $2n$ as a sum of two primes grows roughly like $n / (\ln n)^2$ by the prime number theorem, so large even numbers should have many Goldbach representations.

Yet no proof exists. The difficulty is that the Goldbach conjecture is a \emph{additive} statement about prime numbers, while our best tools for studying primes (the sieve methods, the circle method, the theory of $L$-functions) are \emph{multiplicative} in nature. The parity problem in sieve theory---the significant difficulty that sieves cannot distinguish between numbers with an even or odd number of prime factors---is the core obstruction.

Significant partial results have been achieved:
\begin{itemize}
    \item Vinogradov (1937): Every sufficiently large \emph{odd} integer is the sum of three primes.
    \item Chen (1973): Every sufficiently large even integer is the sum of a prime and a number that is either prime or the product of two primes (``$1 + 2$'').
    \item Helfgott (2013): The \emph{weak} (ternary) Goldbach conjecture is true for all odd integers greater than 5.
    \item Montgomery and Vaughan (2013): The set of even integers that are \emph{not} the sum of two primes has density zero.
\end{itemize}

The Goldbach conjecture remains one of the significant open problems in number theory. Its resolution will likely require fundamentally new ideas about the additive structure of prime numbers.

\subsection{Other Open Problems with Hilbert Connections}

While the Goldbach conjecture is the most prominent survivor of Hilbert's list, other problems have spun off open questions that remain active:

\begin{itemize}
    \item \textbf{The Hodge conjecture} (not one of Hilbert's, but inspired by Problem 16): One of the Clay Millennium Problems, it asks whether certain topological cycles on algebraic varieties are algebraic. It remains completely open.
    \item \textbf{The Langlands program}: Grew from Hilbert's ninth problem and remains the grand unifying vision of number theory. Many of its central conjectures are open.
    \item \textbf{Hilbert's fifth problem in modern guise}: While the classical problem is solved, its generalizations---to infinite-dimensional groups, to approximate groups, and to topological groups with additional structure---are areas of active research.
    \item \textbf{Hilbert's sixteenth problem in modern guise}: The topology of real algebraic varieties, the study of limit cycles of polynomial vector fields (the second part of Problem 16), and connections to Hilbert's eighteenth problem (sphere packing) continue to generate significant mathematics.
\end{itemize}

% ============================================================
\section{The Spirit of Hilbert}
% ============================================================

It is impossible to understand Hilbert's legacy without understanding the \emph{attitude} behind it. Hilbert was not merely a mathematician who solved problems. He was a mathematician who believed in the solvability of problems---who believed that every well-posed mathematical question has an answer, and that human ingenuity will eventually find it.

\subsection{Optimism as a Methodology}

Hilbert's optimism was not naive. It was methodological. He believed that the right way to make progress in mathematics was to \emph{state problems clearly}---to insist on precision and rigor---and then to attack them with energy and imagination. His famous declaration, ``Wir m\"ussen wissen. Wir werden wissen.'' (``We must know. We will know.'') was not a prophecy. It was a \emph{program}: an assertion that the pursuit of knowledge is a moral imperative, and that the obstacles we face are temporary.

This optimism was vindicated many times. When Hilbert posed his problems in 1900, many mathematicians thought that questions about the continuum, about Diophantine equations, and about the foundations of mathematics were too hard---perhaps unanswerable. Within decades, all of them had been addressed, and the answers had enriched mathematics beyond what anyone had imagined.

But Hilbert's optimism also had its limits. G\"odel's theorem showed that there are questions that cannot be answered within any given formal system---that the dream of a complete and consistent foundation for mathematics was, in its original form, impossible. This was not a failure of optimism but a \emph{refinement} of it: it told us that the mathematical universe is more extensive than any single framework can capture.

\subsection{The Culture of Problem-Solving}

Hilbert's list created a \emph{culture} of problem-solving in mathematics. Before 1900, mathematicians communicated results; after 1900, they also communicated \emph{problems}. The idea that a list of open questions could organize an entire field's research agenda became standard. The Fields Medal, first awarded in 1936, is given partly for solving important problems---a direct descendant of Hilbert's vision.

This culture has its dangers. There is always a risk that mathematicians will focus on fashionable problems at the expense of unfashionable ones, or that the prestige attached to solving a famous problem will distort the direction of research. Hilbert himself was aware of this danger---he warned against following ``the moods and fashions of the day'' and urged mathematicians to seek out problems that were ``significant'' rather than merely ``difficult.''

\subsection{Hilbert's Influence on Mathematical Style}

Hilbert's influence extended beyond specific problems to the \emph{style} of mathematics itself. He insisted on complete rigor---every theorem should have a complete proof, with no gaps and no appeals to intuition. This standard, which we now take for granted, was not universal in 1900. Hilbert's \emph{Foundations of Geometry} (1899) was a model of axiomatic precision, and his work on integral equations (1904--1910) established the framework of Hilbert spaces that underpins all of modern functional analysis.

The idea that mathematics should be \emph{clear}---that a proof should be a chain of logical deductions accessible to any competent reader---is Hilbert's gift to the mathematical community as much as any specific theorem.

% ============================================================
\section{New Problems for the 21st Century}
% ============================================================

Hilbert's problems were not the last word. Mathematics continues to generate new questions of significant depth and elegance, and the tradition of posing challenging problems continues.

\subsection{The Clay Millennium Problems}

In 2000, the Clay Mathematics Institute announced seven ``Millennium Prize Problems,'' each with a million-dollar prize for its solution. The problems were chosen to span the major areas of mathematics and to represent the most significant open questions of our time.

\begin{center}
\small
\renewcommand{\arraystretch}{1.2}
\begin{tabular}{cll}
\toprule
\textbf{\#} & \textbf{Problem} & \textbf{Status} \\
\midrule
1 & Birch and Swinnerton-Dyer Conjecture & Open \\
2 & Hodge Conjecture & Open \\
3 & Navier--Stokes Existence and Smoothness & Open \\
4 & P vs. NP & Open \\
5 & Poincar\'e Conjecture & Solved (Perelman, 2003) \\
6 & Riemann Hypothesis & Open \\
7 & Yang--Mills Existence and Mass Gap & Open \\
\bottomrule
\end{tabular}
\end{center}

The most striking aspect of this list is how different it is from Hilbert's. Hilbert's problems were rooted in the mathematics of 1900---classical analysis, number theory, algebraic geometry, the foundations of mathematics. The Millennium Problems reflect the mathematics of 2000---computational complexity, gauge theory, partial differential equations, the geometry of manifolds. The mathematical universe has expanded significantly in a century.

Yet there are continuities. The Riemann Hypothesis, arguably the most famous open problem in all of mathematics, is closely related to Hilbert's eighth problem (the Goldbach conjecture and the distribution of primes). The Birch and Swinnerton-Dyer conjecture is a descendant of Hilbert's tenth problem (Diophantine equations) and his ninth problem (reciprocity laws). The Hodge conjecture is a descendant of Hilbert's sixteenth problem (topology of algebraic varieties). The Poincar\'e conjecture, now solved, was in the spirit of Hilbert's fifth problem (the relationship between topology and algebra).

\subsection{The Langlands Program}

If there is a single vision that defines the research frontier of twenty-first-century number theory, it is the \textbf{Langlands program}---an extensive web of conjectures connecting number theory, algebraic geometry, and the theory of automorphic forms. Named after Robert Langlands, who sketched its outlines in a letter to Andr\'e Weil in 1967, the Langlands program is the heir to Hilbert's ninth problem.

The Langlands program asserts, roughly, that every ``motivic'' object in algebraic geometry corresponds to an ``automorphic'' object in the theory of $L$-functions, and that this correspondence preserves arithmetic information. The most celebrated recent progress is the work of Andrew Wiles (1995) on Fermat's Last Theorem, which proved a special case of the modularity conjecture (itself a case of Langlands). The proof of the Sato--Tate conjecture by Richard Taylor and others, the work of Vincent Lafforgue on the automorphic side, and the recent breakthroughs in $p$-adic Langlands theory all suggest that the Langlands program, while far from complete, is a viable roadmap for the future.

\subsection{Other Grand Challenges}

Beyond the Millennium Problems and the Langlands program, mathematics continues to generate questions of significant depth:

\begin{itemize}
    \item \textbf{The abc conjecture} (now a theorem, per Mochizuki's controversial work): A significant relationship between the additive and multiplicative structure of integers, with consequences for Diophantine equations.
    \item \textbf{The Jacobian conjecture}: Every polynomial map $\mathbb{C}^n \to \mathbb{C}^n$ with everywhere non-zero Jacobian determinant is bijective. Simple to state, possibly impossible to prove.
    \item \textbf{Smooth versus combinatorial 4-manifolds}: In dimension 4, exotic smooth structures exist that have no analog in other dimensions. Understanding these is one of the most significant challenges in topology.
    \item \textbf{Turbulence}: The mathematical theory of the Navier--Stokes equations in three dimensions remains one of the great unsolved problems of mathematical physics, with connections to the Millennium Problem.
    \item \textbf{Quantum computing and quantum gravity}: The link between complexity and physics---``Why is the universe computationally tractable?''---is a significant question Hilbert might have posed today.
\end{itemize}

% ============================================================
\section{Advice for the Reader}
% ============================================================

We began this book with Hilbert's problems as historical objects---questions posed in 1900, studied for over a century, and in some cases still open. We end it with a different perspective: these problems are not historical artifacts. They are \emph{active questions}, and the tradition they represent is one that you, the reader, can participate in.

\subsection{How to Approach Open Problems}

The experience of studying Hilbert's problems suggests several principles for approaching open questions in mathematics:

\begin{enumerate}
    \item \textbf{Start with the special case.} Nearly every major theorem begins with a special case. The weak Goldbach conjecture was proved before the full Goldbach conjecture. The modularity theorem was proved for semistable elliptic curves before the general case. Special cases build intuition, develop techniques, and reveal the structure of the problem.

    \item \textbf{Look for the right framework.} Many of Hilbert's problems were solved not by attacking them directly but by creating new mathematical frameworks---Sobolev spaces for boundary value problems, forcing for set theory, class field theory for reciprocity laws. The right language can make a hard problem tractable.

    \item \textbf{Connect to other areas.} The most significant advances often come from unexpected connections. Hilbert's tenth problem connected number theory to computability theory. Riemann--Hilbert connected differential equations to algebraic geometry. The Langlands program connects number theory with representation theory. If a problem resists direct attack, look for connections to other fields.

    \item \textbf{Be patient.} Hilbert posed his problems in 1900. The Goldbach conjecture is still open in 2026. The Riemann Hypothesis, posed in 1859, remains unsolved. Mathematics operates on a timescale that is longer than a human career. The problems you work on may not be solved in your lifetime---but the tools you develop may help someone else solve them.

    \item \textbf{Don't be afraid of failure.} Almost every major theorem in this book was preceded by failed attempts. Weierstrass's counterexample destroyed the Dirichlet principle---and Hilbert's rescue of it led to the direct method, one of the most important techniques in analysis. G\"odel's theorem was a ``failure'' of Hilbert's program---and it opened up a universe of mathematical logic that Hilbert could not have imagined. Failure is not the opposite of progress; it is a form of progress.
\end{enumerate}

\subsection{The Importance of Persistence}

The story of Hilbert's problems is, above all, a story of persistence. Dehn spent years developing his invariant to resolve Problem 3. G\"odel worked for years on the incompleteness theorems. Cohen developed forcing over a period of intense concentration. Wiles spent seven years in secret work on Fermat's Last Theorem.

Mathematics rewards persistence---but only persistence in the right direction. Hilbert warned against ``the bane of premature specialization'' and urged mathematicians to keep their eyes on the big picture. The most successful problem-solvers are those who combine significant technical skill with a sense of where the important questions are.

\subsection{A Final Word}

Hilbert's problems are not just a list of questions. They are a declaration of faith in the power of human reason. In 1900, Hilbert stood before the International Congress of Mathematicians and declared that the mathematical universe was knowable---that its secrets could be uncovered by patient, rigorous, imaginative work. A century later, we know that he was right about some things and wrong about others. But the faith itself has been vindicated. Mathematics has grown beyond anything Hilbert could have imagined, and the problems he posed---whether solved, unsolved, or transformed beyond recognition---have been at the heart of that growth.

The problems of the twenty-first century are out there, waiting to be posed and to be solved. The mathematical universe is more extensive, richer, and more complex than Hilbert knew. The veil he spoke of has not been lifted---but each theorem, each conjecture, each new idea pushes it back a little further.

\begin{quote}
\textit{``The opinion seems to be gaining ground that this general outlook is very considerably mistaken, and that the best clue to the understanding of the structure of a whole science is to be found in a careful study of some one of its great problems.''}
\par\hfill ---J.\ J.\ Sylvester, 1869
\end{quote}

% ============================================================
\section*{Summary Table: All 23 Problems and Their Status}
% ============================================================

The following table summarizes the current status of all twenty-three of Hilbert's problems.

\begin{center}
\footnotesize
\setlength{\tabcolsep}{3.5pt}
\renewcommand{\arraystretch}{1.15}
\begin{tabular}{c >{\raggedright\arraybackslash}p{3.8cm} >{\raggedright\arraybackslash}p{1.6cm} >{\raggedright\arraybackslash}p{3.8cm}}
\toprule
\textbf{\#} & \textbf{Problem} & \textbf{Status} & \textbf{Key Figures} \\
\midrule
1 & Continuum Hypothesis & Independent & G\"odel, Cohen \\
2 & Consistency of Arithmetic & Resolved (neg.) & G\"odel \\
3 & Equality of Volumes (Tetrahedra) & Solved & Dehn (1901) \\
4 & Shortest Path = Straight Line (Geodesics) & Solved & Hilbert (1904) \\
5 & Continuous Groups = Manifolds (Lie Groups) & Solved & Cartan, von Neumann \\
6 & Axioms of Physics & Partially Solved & von Neumann, Kolmogorov \\
7 & Irrationality of $e^\pi$ & Solved & Gelfond, Schneider (1934) \\
8 & Goldbach Conjecture & Open & Helfgott (weak case) \\
9 & Reciprocity Laws (General) & Solved & Artin, Hasse, Takagi \\
10 & Decision Problem for Dioph. Equations & Solved (neg.) & Davis, Putnam, Robinson, Matiyasevich \\
11 & Quadratic Forms & Partially Solved & Hasse, Minkowski \\
12 & Kronecker's Theorem (Abelian Ext.) & Solved & Takagi, Artin \\
13 & Solving 7th Degree by Radicals & Solved & Kronecker, Hilbert \\
14 & Constructive Algebraic Invariants & Partially Solved & Noether, Kolmogorov \\
15 & Schubert's Enumerative Calculus & Partially Solved & Schubert, Fulton, MacPherson \\
16 & Topology of Algebraic Varieties & Partially Solved & Hodge, Thom, Milnor \\
17 & Representation by Quadratic Forms & Solved & Hasse, Minkowski \\
18 & Tessellation of Space (Keller) & Solved & Debrunner (1965) \\
19 & Regularity of Solutions to PDEs & Partially Solved & Schauder, De Giorgi, Nash \\
20 & Boundary Value Problems & Partially Solved & Hilbert, Sobolev, Lax--Milgram \\
21 & Monodromy Problem & Solved & Deligne, Grothendieck, Riemann--Hilbert \\
22 & Uniformization Theorem & Solved & Poincar\'e, Koebe \\
23 & Calculus of Variations & Partially Solved & Hilbert, Tonelli, Morrey \\
\bottomrule
\end{tabular}
\end{center}

\vspace{1em}

\noindent In total: \textbf{13 problems solved}, \textbf{8 partially solved}, \textbf{1 independent/undecidable}, and \textbf{1 unsolved} (Goldbach). The partially solved problems are, in many ways, the most interesting: they are not failures but invitations---ongoing research programs that continue to generate new mathematics.

% ============================================================
\section*{Exercises}
% ============================================================

\begin{enumerate}

    \item \textbf{The Most Impactful Problem.}
    After reading through the chapters and the retrospective, write a one-page essay answering the following question: \emph{Which of Hilbert's 23 problems has had the greatest impact on modern mathematics, and why?} Consider not only whether the problem was solved, but how the attempts to solve it changed mathematics.

    \item \textbf{Hilbert's Program and G\"odel.}
    Hilbert's second problem asked for a proof of the consistency of arithmetic using finitary methods. G\"odel's incompleteness theorems showed that this is impossible.
    \begin{enumerate}
        \item Explain, in your own words, what G\"odel's first incompleteness theorem says and why it implies that Hilbert's program (in its original form) cannot be carried out.
        \item Discuss: does G\"odel's theorem represent a ``failure'' of Hilbert's vision, or does it represent a \emph{success} of the rigorous methods Hilbert championed?
    \end{enumerate}

    \item \textbf{Comparing Problem Lists.}
    Compare Hilbert's 23 problems (1900) with the Clay Millennium Problems (2000). Write an essay addressing:
    \begin{enumerate}
        \item How do the two lists differ in scope, style, and ambition?
        \item Which areas of mathematics are represented in each list, and which are absent?
        \item What does the comparison tell us about how mathematics has changed in the twentieth century?
    \end{enumerate}

    \item \textbf{The Culture of Problem-Solving.}
    Hilbert's problems helped create a culture of problem-solving in mathematics. Write a one-page essay discussing the \emph{benefits and dangers} of this culture. Consider:
    \begin{enumerate}
        \item Benefits: Does having a shared list of important problems help mathematics progress? How?
        \item Dangers: Does focusing on famous problems distract from less glamorous but equally important work? Can the pursuit of prizes distort research directions?
        \item Are there alternative ways to organize mathematical research that might complement the problem-solving culture?
    \end{enumerate}

    \item \textbf{Hilbert's Legacy in Your Field.}
    Choose one area of mathematics that interests you (e.g., number theory, topology, analysis, combinatorics, algebra, logic, applied mathematics). Write a one-page essay describing:
    \begin{enumerate}
        \item Which of Hilbert's problems, if any, are relevant to this area today?
        \item What are the most important open problems in this area as of 2026?
        \item What new techniques or ideas have emerged in this area that Hilbert could not have anticipated?
    \end{enumerate}

    \item[$\star$] \textbf{Designing a Problem List for 2050.}
    If you were to design a list of 23 important open problems for mathematics in the year 2050, what would you include? Write a one-page proposal for such a list. Your proposal should:
    \begin{enumerate}
        \item Name at least five specific open problems (they need not be original).
        \item Explain why each problem is important.
        \item Discuss what areas of mathematics your list emphasizes, and what areas it omits. Is your list as balanced as Hilbert's?
        \item Reflect on whether your list is a list of \emph{problems} (specific questions with definite answers) or a list of \emph{programs} (broad research directions). What are the advantages and disadvantages of each approach?
    \end{enumerate}

\end{enumerate}
